\documentclass[11pt]{article}
\usepackage[letterpaper,margin=0.88in]{geometry}
\usepackage[T1]{fontenc}
\usepackage{lmodern,microtype,amsmath,amssymb,amsthm,booktabs,array,needspace}
\usepackage[hidelinks]{hyperref}
\hypersetup{pdfauthor={Tynan Daly}}
\usepackage{fancyhdr}
\pagestyle{fancy}
\fancyhf{}
\fancyhead[L]{\small Recurrence and stationarity on cycles}
\fancyhead[R]{\small Limited pass-and-swap queues}
\fancyfoot[C]{\thepage}
\setlength{\headheight}{14pt}
\setlength{\parindent}{1.2em}
\setlength{\parskip}{3pt}
\setlength{\emergencystretch}{1.5em}
\newtheorem{theorem}{Theorem}
\newtheorem{lemma}[theorem]{Lemma}
\newtheorem{proposition}[theorem]{Proposition}
\newtheorem{corollary}[theorem]{Corollary}
\theoremstyle{remark}
\newtheorem{remark}[theorem]{Remark}
\newcommand{\rev}{\operatorname{rev}}
\newcommand{\B}{\mathcal B_w}
\newcommand{\W}{W_\theta}
\newcommand{\F}{F_\theta}
\newcommand{\Qw}{Q_{w,\theta}}
\newcommand{\Qi}{Q_{\infty,\theta}}
\newcommand{\emptyq}{\varnothing}

\newcommand{\T}{\mathcal T_{n,w}}
\title{Recurrence and Product-Form Stationarity\\in Swap-Limited Queues on Cycles}
\author{Tynan Daly}
\date{September 7, 2026}
\begin{document}
\maketitle
\begin{abstract}
We classify the recurrent configurations of two pass-and-swap queues in closed tandem, with one distinct job per vertex of a cycle. For a swap limit $w\ge1$ and positive service at every occupied position, each orientation whose circular runs have length at most $w$ gives one closed communicating class. The only other recurrent configurations have exactly one run of length $w+1$ and all other runs of length $w$. They form a unique exceptional class, which exists exactly when $n=2kw+1$ for an integer $k\ge1$. For $w\ge2$, this yields a sharp boundary: the canonical OI weights are stationary on every closed class for every positive OI allocation if and only if $n\not\equiv1\pmod{2w}$. At exceptional lengths, complementary constructions give exact canonical balance defects: a distinguished class rate on $C_{2w+1}$, and unit position rates on every larger exceptional cycle. Integer stationary certificates on $C_5$ additionally exclude every queue-wise factorization, both at a fixed budget and when an independent two-state Markov process switches between limited and unlimited swaps. The latter gives a counterexample to the proposed Markov-modulated product form. The proofs connect actual completions to orientation moves and run dynamics; the arbitrary-size classification and sharp boundary are verified in Lean, with exact finite computations retained separately.
\end{abstract}

\section{Introduction and main results}\label{sec:intro}
Limiting a pass-and-swap replacement chain can change both the recurrent configurations and their stationary probabilities. The central question is which configurations survive indefinitely once a finite swap budget is imposed. For cycle swapping graphs, we give a complete answer and use it to determine precisely when the canonical product form is guaranteed for every strictly positive OI allocation.

Pass-and-swap queues model scheduling and load-distribution protocols in machine clusters where compatibility constraints determine which jobs each machine can serve~\cite[pp.~277--278]{CD}. Their replacement chains describe how a service completion can move jobs through a sequence of compatible positions. A product-form stationary law expresses a configuration's probability through explicit queue-wise factors and one normalizing constant, enabling exact performance calculations without solving a separate balance equation for every state.

Comte and Dorsman~\cite{CD} introduced pass-and-swap (P\&S) queues as an extension of order-independent (OI) queues. Dorsman and Gardner~\cite{DG} proved that their canonical stationary law survives a swap limit when the initial placement order is short, and conjectured transience of taller configurations under a partiteness hypothesis. Our classification identifies all short classes and the unique possible tall class. It places the earlier balanced family $\mathcal B_w$ in the larger family $\mathcal T_{n,w}$, and separates the structural recurrence question from the allocation-dependent balance question.

\paragraph{Queues and completions.}
A configuration is $s=(c,d)$, with each queue written from head to tail. There is one distinct job with each label in $\{0,\ldots,n-1\}$, and the swapping graph is the simple cycle $C_n$, $n\ge3$, with edges $\{i,i+1\pmod n\}$. A completion may occur at any occupied position. Its job replaces the first later compatible job, which continues the same scan strictly farther back. After $w$ replacements, or when no later compatible job remains, the carried job departs and joins the other queue's tail. The initiating completion does not consume a replacement. These are the rules of~\cite[\S2.3, pp.~211--212; \S3.1, p.~215; \S5, p.~234]{DG}, with the original position-index recursion in~\cite[\S3.1, pp.~284--285]{CD}. Unless explicitly stated otherwise, completion rates are finite, time-homogeneous, and strictly positive at every occupied position. The recurrence theorem permits dependence on the entire configuration and does not require OI service.

\paragraph{Orientations and runs.}
Orient each cycle edge from its earlier to its later endpoint in
\begin{equation}\label{eq:order}
 L(c,d)=c\Vert\rev(d).
\end{equation}
This is the placement orientation for distinct labels; its transitive closure is the placement order of~\cite[\S3.1, p.~216]{DG}. Encode it by the circular word $(b_0,\ldots,b_{n-1})$, where $b_i=1$ means $i\to i+1$. The two constant words are directed cycles and cannot arise from a linear order. Every other word is an acyclic orientation and is realized by some configuration. A \emph{run} is a maximal circular block of equal bits. There are an even number $2k$ of runs, and the placement height $h(s)$, measured in edges, is their maximum length.

\begin{theorem}[Complete recurrent-class classification]\label{thm:classification}
Let $n\ge3$ and $w\ge1$. Under strictly positive service at every occupied position, a configuration is recurrent if and only if its circular run lengths satisfy one of the following conditions:
\begin{enumerate}
\item Every run has length at most $w$.
\item Exactly one run has length $w+1$, and every other run has length $w$.
\end{enumerate}
Each orientation in the first case gives one closed communicating class, consisting of all its linear extensions and all their queue cuts. A class of the second kind exists if and only if
\begin{equation}\label{eq:exceptional}
 n=2kw+1\qquad\text{for an integer }k\ge1.
\end{equation}
When it exists, it is unique: $\T$ consists of every configuration with this exceptional run pattern. It contains exactly $2n$ orientations, with all their linear extensions and cuts, and every state in it has height $w+1$. There are no other recurrent states. From any initial configuration, the process almost surely eventually remains in one of these classes.
\end{theorem}

The exceptional set is closed under \emph{every} position completion. Thus closure, and existence of a recurrent state in this nonempty finite set, persist under arbitrary nonnegative allocations. General irreducibility and the exhaustive recurrence classification require the positivity stated in Theorem~\ref{thm:classification}. The stronger head-only communication theorem for the five-job class is retained separately in Appendix~\ref{app:head}.

\paragraph{Canonical weights.}
For product-form statements, let $\mu$ and $\nu$ be OI total capacities: they vanish on the empty queue, depend only on the multiset of labels, and have strictly positive prefix increments on every valid nonempty prefix. The position-$j$ completion rate in the first queue is
\[
 \mu(c_1,\ldots,c_j)-\mu(c_1,\ldots,c_{j-1}),
\]
and similarly for $\nu$. The two allocations may differ. On a closed class $\mathcal C$ define
\begin{equation}\label{eq:canonical}
 W(c,d)=\prod_{j=1}^{|c|}\frac1{\mu(c_1,\ldots,c_j)}
         \prod_{j=1}^{|d|}\frac1{\nu(d_1,\ldots,d_j)},\qquad
 \widehat\pi_{\mathcal C}(s)=\frac{W(s)}{Z_{\mathcal C}},\quad
 Z_{\mathcal C}=\sum_{s\in\mathcal C}W(s)>0.
\end{equation}
Empty products equal one. These are the unnumbered prefix factors immediately below Eq.~(25) of~\cite[p.~240]{DG}. Stationarity means $\widehat\pi_{\mathcal C}Q=0$ for the original continuous-time generator on $\mathcal C$; generators act on row vectors.

\Needspace{9\baselineskip}
\begin{theorem}[Sharp universal canonical boundary]\label{thm:sharp}
For $n\ge3$ and $w\ge2$, the following are equivalent:
\begin{enumerate}
\item $n\bmod(2w)\ne1$.
\item For every pair of strictly positive OI allocations in the two queues, the normalized canonical weights~\eqref{eq:canonical} are stationary on every closed communicating class.
\end{enumerate}
\end{theorem}

This is universal validity over allocations. At an exceptional length it asserts the existence of an allocation violating canonical stationarity, not failure for every allocation. It is not a criterion for stationarity of an individually specified allocation. The sufficiency direction holds already for $w\ge1$ by Corollary~\ref{cor:positive}\footnote{Formalized as \texttt{ClosedClass.safe\_normalized\_oi}; the sharpness equivalences require $w\ge2$.}: every nonexceptional cycle, in particular every even cycle, has only short recurrent classes. A law started outside those classes may converge to a mixture of classwise stationary laws; one globally normalized $W$ is not asserted.

\paragraph{Relation to Conjecture 1.}
Conjecture~1 of~\cite[p.~241]{DG} proposes transience of tall states on a $(w+1)$-partite graph. Every cycle is three-colorable, hence $(w+1)$-partite for $w\ge2$, where partiteness means a coloring with at most $w+1$ colors, as in~\cite[Lemma~3, pp.~252--253]{DG}. The recurrence results give the following boundary, with different service assumptions in its two directions.

\begin{corollary}[Cycle boundary for Conjecture 1]\label{cor:conjecture}
For $n\ge3$ and $w\ge2$, with one job per vertex of $C_n$, Conjecture~1 fails for every admissible OI allocation, allowing zero non-head rates, whenever $n\equiv1\pmod{2w}$. For every other $n$, it holds under strictly positive service at every occupied position.
\end{corollary}
\begin{proof}
At an exceptional length, $\T$ is nonempty and closed under every position completion. Under any admissible OI allocation it therefore contains a recurrent state, necessarily of height $w+1$. At every other length, Theorem~\ref{thm:classification} makes all tall states transient when position rates are strictly positive.
\end{proof}

Removing positivity from the validity direction remains open. Dorsman and Gardner explicitly note that their OI conditions guarantee positive head rates only~\cite[Appendix~C.2, p.~254]{DG}. Our finite checks in Section~\ref{sec:finite} find no tall recurrent states on the tested nonexceptional cycles under head-only service either, suggesting that this hypothesis may be removable. At $w=1$, the classification confirms Conjecture~1 on every even cycle under positive position rates; Corollary~\ref{cor:even-one} identifies its two closed classes. Odd cycles fail the bipartite hypothesis. The arbitrary-graph one-swap bipartite problem and removal of the positive-position assumption remain unresolved. Theorem~7 and Eq.~(25) of~\cite[p.~240]{DG}, which concern short initial placement orders, are unaffected. Counterexample~2 of~\cite[p.~239]{DG} gives product-form failure outside the relevant partite regime.

The proof proceeds from actual completions to exact orientation moves (Section~\ref{sec:quotient}), then from circular runs to terminal classes (Section~\ref{sec:runs}). Short-class balance proves the positive direction (Section~\ref{sec:positive}); complementary exceptional-family calculations prove sharpness (Section~\ref{sec:failures}). The five-job exclusion of all queue-wise factors also extends to an independent two-state modulation switching between budget two and unlimited service. Section~\ref{sec:modulated} gives an exact counterexample to Conjecture~2 of~\cite[p.~245]{DG}, allowing even the candidate queue factors to depend on the environment state. The fixed-budget certificate, this modulation result, a short queue-length consequence, and finite verification follow in Sections~\ref{sec:certificate}--\ref{sec:finite}.

\section{From queue completions to orientation moves}\label{sec:quotient}
We first discard the queue cut and the choice of linear extension without discarding reachability information. This is a quotient of an event graph: a quotient edge exists when some actual completion realizes it. No claim that the orientation projection is itself Markov is needed.

\begin{lemma}[Communication within an orientation]\label{lem:fiber}
All configurations inducing a fixed acyclic orientation communicate by completions making zero replacements.
\end{lemma}
\begin{proof}
Tail completions move the cut in $L$ one place in either direction, preserving $L$. For adjacent incomparable labels $a,b$ in a linear extension $PabT$, move the cut immediately after the pair. Completing $a$ in $(Pab,\rev(T))$ makes no replacement, because incomparable labels cannot share a swapping edge. The result $(Pb,\rev(T)a)$ has placement word $PbaT$. Tail completions restore any desired cut, and the reverse exchange is equally available.

Any two linear extensions are connected by adjacent incomparable exchanges. To see this, move the next label of the desired extension leftwards until its position agrees. Each label it crosses occurs in the reverse order in two linear extensions, so the labels are incomparable. Iterating gives the desired extension. Positive position rates make all the displayed event paths available.
\end{proof}

For a finite event graph partitioned into communicating fibers, quotient paths lift to full configuration paths: move inside the current fiber to a representative realizing the next edge, perform that completion, and continue. Conversely, every full path projects to a quotient walk. Communication and terminal components therefore correspond exactly. The quotient need not have well-defined transition rates. For example, on $C_5$ at $w=2$, $((0,1,2,4,3),\emptyq)$ and $((0,1),(3,4,2))$ induce the same orientation, but only the former admits an orientation-changing completion.

\begin{lemma}[Exact nontrivial quotient moves]\label{lem:local}
For $n\ge3$ and $w\ge1$, a nontrivial orientation move flips one bit $b_i$. It is realizable if and only if the $w$ immediately preceding bits all equal $b_i$, or the $w$ immediately following bits all equal $b_i$, with circular indices. Equivalently, it flips an end edge of a directed path of $w+1$ edges, with queue exchange allowing the two ends. If $w\ge n-1$, no such move exists.
\end{lemma}
\begin{proof}
Consider a completion in the first queue. Its replacement chain follows a directed path in the placement orientation. A carried label crosses no compatible label before its first later replacement target, and replacing that target retains the direction of their edge. Compaction changes no other pair's relative order. Consequently only an edge from the final carried label to a later compatible label can change direction. If the scan stops before exhausting the budget, no such later label exists and the orientation is preserved, including all edges between the queues.

If the scan makes $w$ replacements, the final displaced vertex has its previous chain neighbor behind it. Degree two leaves at most one further neighbor. Its edge reverses exactly when that neighbor is still later in the completing queue; otherwise no orientation changes. A changed edge is preceded by the $w$ directed chain edges. Exchanging the queues reverses the placement word and gives the condition from the other end. This proves necessity directly from the operational scan.

For sufficiency, take a directed path $v_0\to\cdots\to v_{w+1}$. Put all jobs in the first queue in a topological order with $v_1$ the first later neighbor of $v_0$. Such an order exists. If the other neighbor $z$ of $v_0$ is outgoing, no directed path from $z$ can force $z$ before $v_1$: a path entering $v_1$ without $v_0$ would have to use its other incident edge, contrary to $v_1\to v_2$. Thus we may place $v_1$ before $z$. An incoming neighbor is already earlier. After the first replacement, degree two forces the displayed path. Completing $v_0$ with budget $w$ flips its final edge. Queue exchange realizes the other-end move, and Lemma~\ref{lem:fiber} makes these representatives accessible. A directed path of $w+1\ge n$ edges cannot occur in an acyclic orientation on $n$ vertices, proving the last assertion.
\end{proof}

Thus a bit at position $j$ of a run of length $\ell$ can flip precisely when
\begin{equation}\label{eq:flip}
 j\ge w+1\quad\text{or}\quad j\le\ell-w,\qquad 1\le j\le\ell.
\end{equation}
The single-edge conclusion uses the cycle's degree bound. For contrast, on the star with edges $\{0,1\},\{1,2\},\{1,3\}$, the $w=1$ transition $((0,1,2,3),\emptyq)\to((0,2,3),(1))$ reverses both $1\to2$ and $1\to3$.

\section{Circular runs and the recurrent classes}\label{sec:runs}
The exact rule~\eqref{eq:flip} gives a short proof of the classification. A permitted flip never removes an entire run: its donor has length at least $w+1\ge2$. Flipping a boundary bit transfers one edge to the adjacent run and preserves run count. Flipping an interior bit splits a run into three and increases run count by two. Run count is therefore nondecreasing along every quotient path.

If some run has length at least $w+2$, its second bit can flip: at least $w$ equal bits follow it. This is an interior flip and permanently increases run count. The original orientation cannot be in a terminal communicating component. It remains to consider run lengths at most $w+1$, where only the first or last edge of a length-$(w+1)$ run can move.

Call a run \emph{full} if its length is $w$, \emph{overfull} if it is $w+1$, and \emph{deficient} if it is smaller than $w$. A boundary move from an overfull run makes that donor full and increases its neighbor by one. If the neighbor is full, this transports the extra edge one run farther. If it is deficient, the deficit
\[
 D=\sum_i\max\{w-\ell_i,0\}
\]
decreases. This decrease cannot be reversed while run count stays fixed: boundary donors have length greater than $w$ and remain at least $w$, so no boundary move increases $D$. A split can create deficits, but it permanently changes run count. If the neighbor is instead overfull, the transfer creates length $w+2$, allowing an irreversible split.

Suppose a terminal component contains an overfull run and another nonfull run. Transport the extra edge through consecutive full runs until the first other nonfull run is reached. Either a deficit decreases or a split becomes available. Each case produces a reachable orientation that cannot return, a contradiction to terminality. Thus a tall terminal orientation must have exactly one overfull run and every other run full. Since a nonconstant circular binary word has $2k$ runs, their sum is $(w+1)+(2k-1)w=2kw+1$.

Conversely, every short orientation has no nontrivial quotient move. Each lifts, by Lemma~\ref{lem:fiber}, to its entire closed communicating fiber. The single-overfull pattern is closed because its only moves transfer the extra edge to a full neighbor. Parameterize such a word by $(a,b)$, with $a$ the starting edge of its unique long run and $b$ that run's bit. Its two boundary moves are
\[
 (a,b)\longrightarrow(a+w,1-b)\quad\text{or}\quad(a-w,1-b)\pmod n.
\]
When $n=2kw+1$, $\gcd(n,w)=1$ and $n$ is odd. Repeating the first move visits all $2n$ pairs: returning to the original edge requires a multiple of $n$ steps, and returning to its bit requires an even number. Every pair specifies a distinct orientation, since the long run is unique. This proves both the count and communication of the exceptional orientations. Path lifting then proves communication of all their extensions and cuts.

We have exhausted the terminal components of the event graph. Strictly positive rates make its edges exactly the positive-rate transitions. A finite continuous-time chain with these finite rates is nonexplosive; its recurrent states are precisely the terminal communicating classes, and it almost surely eventually stays in one such class. This proves Theorem~\ref{thm:classification}, including its stochastic assertions.

\paragraph{A linear recurrence test.}
Build the ranks in $c\Vert\rev(d)$, read the $n$ edge bits, and scan the circular run lengths. Theorem~\ref{thm:classification} decides recurrence in $O(n)$ bounded-word RAM operations, without enumerating $(n+1)n!$ configurations. The instrumented Lean algorithm has cost at most $33n+18$; array accesses and arithmetic on words large enough for the input quantities count as constant cost. This is not a bit-complexity or execution-time claim. The congruence decides whether a tall class exists, not whether every tall configuration belongs to it.

\section{Canonical stationarity on the short classes}\label{sec:positive}
\begin{corollary}[Positive product form]\label{cor:positive}
Let $n\ge3$, $w\ge1$, and $n\bmod(2w)\ne1$. For every pair of strictly positive OI allocations, the normalized canonical weights are stationary on each recurrent class. Every initial configuration almost surely enters a short-orientation class.
\end{corollary}
\begin{proof}
In a short orientation, every directed path has at most $w$ edges. Thus every replacement chain has at most $w$ replacements, and limited and unlimited dynamics coincide. The orientation fiber is closed and communicating by Lemma~\ref{lem:fiber}. The unlimited tandem theorem, Theorem~3 and Eq.~(8) of~\cite[p.~217]{DG}, gives the canonical law on this fiber; its original version is Theorem~5 and Eq.~(23) of~\cite[p.~298]{CD}. The same mechanism underlies Theorem~7 of~\cite[p.~240]{DG}, which is stated with an additional partiteness hypothesis. The classification shows that these are all recurrent classes at nonexceptional lengths and supplies the asserted absorption. In the formal development the unlimited OI balance identity is proved from complete incoming-event reconstruction and telescoping prefix weights, rather than assumed from the literature.
\end{proof}

\begin{corollary}[The one-swap conjecture on even cycles]\label{cor:even-one}
Let $n\ge4$ be even, with one job per vertex of $C_n$ and swap budget $w=1$. Under strictly positive service at every occupied position, there are exactly two closed communicating classes, corresponding to the two alternating cycle orientations. Each class contains all linear extensions and cuts of its orientation; every other configuration is transient. For every pair of positive OI allocations, the normalized canonical weights give the stationary law on each class. Thus Conjecture~1 holds on every even cycle in this positive-position regime.
\end{corollary}
\begin{proof}
An even $n$ is not congruent to $1$ modulo $2$, so Theorem~\ref{thm:classification} leaves only short orientations. At $w=1$ every circular run then has length one. There are exactly two such circular binary words, the two alternating orientations; both are acyclic and realizable. The classification identifies their full configuration classes and transience outside them, and Corollary~\ref{cor:positive} gives their canonical laws.
\end{proof}

This proves sufficiency in Theorem~\ref{thm:sharp}. The next section proves necessity by giving one failing positive allocation on every exceptional cycle. A nonzero \emph{unnormalized} residual remains nonzero after division by the positive class normalizer.

\section{Canonical failure on every exceptional cycle}\label{sec:failures}
At the first exceptional length, $n=2w+1$, one distinguished class rate gives a uniform calculation. At every larger exceptional length, unit position rates suffice.

\subsection{The first exceptional class: \texorpdfstring{$\mathcal B_w=\mathcal T_{2w+1,w}$}{the balanced family}}\label{sec:balanced}
Fix $w\ge2$. The exceptional pattern has two runs, or equivalently two source-to-sink branches of lengths $w$ and $w+1$. Denote this class by $\B$. This notation is reserved for $n=2w+1$; the classes with $k\ge2$ have multiple sources and sinks. Theorem~\ref{thm:classification} proves recurrence and communication of $\B$ under positive position rates.
\begin{lemma}[Allocation-independent closure]\label{lem:closure}
The set $\B$ is nonempty, has cardinality
\begin{equation}\label{eq:size}
 |\B|=2(2w+1)(2w+2)\binom{2w-1}{w},
\end{equation}
and is closed under completions at every queue position, irrespective of the nonnegative rates assigned to those completions.
\end{lemma}
\begin{proof}
Compare a completion in $c$ with its unlimited version. The unlimited chain follows a directed path and therefore uses at most $w+1$ replacements. If it uses at most $w$, the transitions agree, and orientation preservation~\cite[Lemma~1, p.~217]{DG} proves that the successor remains balanced.

Otherwise, the chain traverses the full longer branch from the unique source $u$ to the unique sink $v$. Both endpoints are in $c$. Since every vertex lies between them in $L(c,d)$, all jobs must be in $c$, with $u$ its head and $v$ its tail. Let $p$ be the predecessor of $v$ on the longer branch. The limited transition ejects $p$, whereas the unlimited transition makes one more replacement and ejects $v$. For the same remaining word $R$, the resulting placement orders are
\[
 L_{\rm limited}=Rvp,\qquad L_{\rm unlimited}=Rpv.
\]
Only the edge $p\to v$ changes orientation. The old longer branch loses one edge and the shorter branch gains the reversed edge; their lengths are again $w,w+1$, now ending at $p$. Exchanging the two queues reverses every oriented edge, so the same argument applies to completions in $d$.

For the count, there are $2(2w+1)$ choices of source and long-branch direction. The $w$ internal vertices of the longer branch interleave with the $w-1$ internal vertices of the shorter branch in $\binom{2w-1}{w}$ ways. The source must be first and the sink last. Each linear extension has $2w+2$ cuts, with the suffix reversed to obtain $d$. This proves~\eqref{eq:size} and nonemptiness.
\end{proof}

Give each position its job's class rate, identically in both queues:
\begin{equation}\label{eq:rates}
 r_0=\theta>0,\qquad r_i=1\quad(1\le i\le2w),\qquad
 \mu(q)=\nu(q)=\sum_{x\in q}r_x.
\end{equation}
These functions are OI: their prefix increments are precisely the positive $r_x$. Define
\begin{equation}\label{eq:weights}
 \F(q)=\prod_{j=1}^{|q|}\frac{1}{\sum_{i=1}^{j}r_{q_i}},\qquad
 \F(\emptyq)=1,\qquad \W(c,d)=\F(c)\F(d).
\end{equation}
Here $W_\theta$ is the specialization of~\eqref{eq:canonical} to these identical additive OI allocations.

\begin{theorem}[Uniform canonical product-form obstruction]\label{thm:balanced}
For every $w\ge2$ and $\theta>0$, $\B$ is a single closed communicating class for~\eqref{eq:rates}. Every state in it has height $w+1$. Set
\[
 d_w=(w,w-1,\ldots,1,w+1,w+2,\ldots,2w),\qquad s_w=((0),d_w).
\]
The limited generator on $\B$ satisfies
\begin{equation}\label{eq:main}
 (\W\Qw)(s_w)=
 \frac{1}{w!\prod_{k=w}^{2w}(\theta+k)}
 -\frac{1}{(2w)!(\theta+2w)}.
\end{equation}
This is positive for $0<\theta<1$ and negative for $\theta>1$. Thus, for every $\theta\ne1$, the unique stationary law on $\B$ is not proportional to $\W$. At $\theta=2$ the defect is
\begin{equation}\label{eq:factorial}
 \boxed{(W_2Q_{w,2})(s_w)=-\frac{w}{(2w+2)!}.}
\end{equation}
\end{theorem}
\subsection{The two-flow proof}\label{sec:balance}
Let $\Qi$ be the unlimited generator on $\B$. Each fixed-orientation set is closed under unlimited P\&S and is irreducible by the zero-replacement argument just given. The unlimited product-form theorem~\cite[Theorem~3, Eq.~(8), p.~217]{DG} gives $\W\Qi=0$ on each such set and thus on their union. Normalizations may differ between orientations; the \emph{unnormalized} balance identity remains valid. Consequently
\begin{equation}\label{eq:subtract}
 \W\Qw=\W(\Qw-\Qi).
\end{equation}
Only events whose destinations change can contribute.

\paragraph{The exchanged destinations.}
Define
\begin{align*}
 x_w&=(w,w-1,\ldots,1,0,w+1,w+2,\ldots,2w),\\
 y_w&=(w,w+1,w-1,\ldots,1,w+2,\ldots,2w,0),\\
 X_w&=(\emptyq,x_w),\qquad Y_w=(\emptyq,y_w),\\
 s'_w&=((2w),(w,w-1,\ldots,1,w+1,\ldots,2w-1,0)).
\end{align*}
Empty ranges in these expressions are omitted. All four states $X_w,Y_w,s_w,s'_w$ are balanced. The head completions at $X_w,Y_w$ have rate one, because the initiating label is $w$, not $0$. Their destinations are
\begin{center}
\begin{tabular}{lcc}
\toprule
Event & Limit $w$ & Unlimited\\
\midrule
Head of $d$ at $X_w$ & $s_w$ & $s'_w$\\
Head of $d$ at $Y_w$ & $s'_w$ & $s_w$\\
\bottomrule
\end{tabular}
\end{center}
Indeed, the unrestricted replacement chains are respectively
\[
 w,w-1,\ldots,1,0,2w
 \quad\hbox{and}\quad
 w,w+1,w+2,\ldots,2w,0.
\]
Each uses $w+1$ replacements; suppressing its last replacement gives the limited destination.

\Needspace{9\baselineskip}
\begin{lemma}[Completeness]\label{lem:two}
Among events whose limited and unlimited destinations differ, the only gained incoming event at $s_w$ is the head completion at $X_w$, and the only lost incoming event is the head completion at $Y_w$.
\end{lemma}
\begin{proof}
The closure proof shows that every changed event from $\B$ starts at a queue containing all jobs, completes its head, and follows an unrestricted chain of length $w+1$. To arrive at $s_w=((0),d_w)$, that queue must be $d$.

For a limited arrival, the ejected penultimate vertex is $0$ and the unchanged final vertex is the last entry of $d_w$, namely $2w$. The unique simple cycle path of length $w+1$ with final edge $0\to2w$ is $w,w-1,\ldots,0,2w$. Reversing its replacements in the target word uniquely reconstructs $x_w$.

For an unlimited arrival, the ejected final vertex is $0$ and its predecessor is the last entry $2w$ of the output word. The chain is therefore $w,w+1,\ldots,2w,0$, which uniquely reconstructs $y_w$. In either reconstruction, each retained chain label occupies the former position of the next chain label. Replacing these labels back and prepending the removed head recovers exactly one source word. There are no other changed incoming events.
\end{proof}

Imposing the limit retargets events but does not change their rates. Every event changes the queue lengths, so no self-transition is created and the generator diagonals agree. Combining~\eqref{eq:subtract} with Lemma~\ref{lem:two} gives
\begin{equation}\label{eq:two}
 (\W\Qw)(s_w)=\W(X_w)-\W(Y_w).
\end{equation}
The calculation has two terms for every $w$.

\paragraph{Evaluation and sign.}
In $x_w$, job $0$ appears at position $w+1$. Its first $w$ prefix totals are $1,\ldots,w$; the remaining totals are $\theta+w,\ldots,\theta+2w$. In $y_w$, job $0$ appears last. Hence
\[
 \W(X_w)=\frac{1}{w!\prod_{k=w}^{2w}(\theta+k)},\qquad
 \W(Y_w)=\frac{1}{(2w)!(\theta+2w)},
\]
proving~\eqref{eq:main}. To determine the sign, put $P_w(\theta)=\prod_{k=w}^{2w-1}(\theta+k)$. This is strictly increasing for $\theta>0$, and $P_w(1)=(2w)!/w!$. Thus
\[
 (\W\Qw)(s_w)=\frac{1}{\theta+2w}
 \left(\frac{1}{w!P_w(\theta)}-\frac{1}{(2w)!}\right)
\]
is positive below $1$ and negative above $1$. Division by the positive normalizing constant $Z_{w,\theta}=\sum_{s\in\B}\W(s)$ cannot remove this defect. Finally,
\[
 W_2(X_w)=\frac{w+1}{(2w+2)!},\qquad
 W_2(Y_w)=\frac{2w+1}{(2w+2)!},
\]
which proves~\eqref{eq:factorial} and completes Theorem~\ref{thm:balanced}.


\subsection{The larger exceptional classes: unit position rates}\label{sec:larger}
Fix $w\ge1$, $k\ge2$, $n=2kw+1$, and $m=2w(k-1)=n-2w-1$. Give every occupied position rate one in both queues, so $\mu(q)=\nu(q)=|q|$ and
\[
 W(c,d)=\frac1{|c|!\,|d|!}.
\]
Choose the orientation whose first clockwise run has length $w+1$, followed by alternating runs of length $w$. Use its realizing topological order from the orientation-realization construction, restrict it to the labels $\{2w+2,\ldots,n-w-1\}$, and call the resulting word $H$. This fixes the interior order used by the formal theorem; empty intervals are omitted. Put $R=\rev(H)$ and
\begin{align*}
 P&=(n-w,n-w+1,\ldots,n-1,R),& |P|&=m-1,\\
 d&=(P,w+1,w+2,\ldots,2w+1,w,w-1,\ldots,1),& s&=((0),d).
\end{align*}
The placement word is
\[
 (0,1,\ldots,w,\,2w+1,2w,\ldots,w+1,\,H,\,n-1,n-2,\ldots,n-w).
\]
It realizes the specified orientation: $2w+1$ is a source at the left boundary of the interior, and $n-w$ is a sink at its right boundary. Hence $s\in\T$.

\begin{theorem}[Unit-rate exceptional-family defect]\label{thm:unit}
With these choices and the generator restricted to $\T$,
\begin{equation}\label{eq:unit}
 (WQ)(s)=\frac{1-2w(k-1)}{n!}<0.
\end{equation}
In particular, the normalized canonical weights are not stationary on $\T$.
\end{theorem}
\begin{proof}
Unlimited dynamics preserve each orientation, by~\cite[Lemma~1, p.~217]{DG}, so their canonical weights balance on the union of the exceptional fibers. Subtract that unlimited balance equation from the limited one. Diagonals agree because truncation changes destinations but not initiating rates, and every event changes a queue length. A changed arrival into $s$ must come from a configuration with all jobs in the second queue: a first-queue predecessor has only two jobs and cannot exceed the budget $w\ge1$.

There is exactly one gained incoming event, the second-queue head completion from
\[
 X=(\emptyq,(n-w,n-w+1,\ldots,n-1,0,R,w+1,\ldots,2w+1,w,\ldots,1)).
\]
Its unrestricted chain is $n-w,n-w+1,\ldots,n-1,0,1$. The limited event ejects $0$ and enters $s$; the unrestricted event ejects $1$. To check completeness, invert the scan at a changed limited arrival. The neighbors of $0$ in the target queue are $n-1$, at position $w$, and $1$, at its tail. Reversing the length-$w$ replacement chain reconstructs exactly the prefix $n-w,\ldots,n-1,0$. Its retained initial label is the head of $d$, forcing the initiating position to be the head. This recovers $X$ uniquely. Its orientation differs from the target at $\{0,1\}$ and has the exceptional pattern, so this gained event is on the stated support.

There are exactly $m$ lost incoming events. Form
\[
 z=(P,w,w+2,\ldots,2w+1,w-1,w-2,\ldots,0).
\]
Insert $w+1$ into any of the $|P|+1=m$ positions before the displayed $w$, and complete that inserted job. The unrestricted chain is $w+1,w,w-1,\ldots,1,0$; it enters $s$, whereas truncation ejects $1$. The prefix $P$ contains neither neighbor $w$ nor $w+2$ of the inserted label, so every insertion works. Conversely, reverse the unrestricted scan from $0$ in $d$. Its chain has more than $w$ replacements exactly for insertion positions before that first block boundary. Reversing its replacements gives precisely these $m$ words. Unlimited orientation preservation places all these sources in the same orientation fiber as $s$.

All full-queue sources have weight $1/n!$ and every event has rate one. All unchanged contributions cancel, leaving $(1-m)/n!$. Since $w\ge1$ and $k\ge2$, $m\ge2$, proving strict negativity.
\end{proof}

\subsection{An explicit unit-rate instance on \texorpdfstring{$C_9$}{C9}}\label{sec:nine}
Take $n=9$, $w=2$, and edges $\{0,1\},\{1,2\},\ldots,\{7,8\},\{8,0\}$. Each position receives rate one in both queues. The target and gained source are
\[
 s=((0),(7,8,6,3,4,5,2,1)),\qquad
 X=(\emptyq,(7,8,0,6,3,4,5,2,1)),
\]
with a head completion at $X$. The four lost source words, all in the second queue with the first empty, are
\begin{center}
\begin{tabular}{lc}
\toprule
Second-queue word & Initiating position\\
\midrule
$(3,7,8,6,2,4,5,1,0)$&1\\
$(7,3,8,6,2,4,5,1,0)$&2\\
$(7,8,3,6,2,4,5,1,0)$&3\\
$(7,8,6,3,2,4,5,1,0)$&4\\
\bottomrule
\end{tabular}
\end{center}
Positions are one-based. Formula~\eqref{eq:unit} specializes to the unnormalized defect
\begin{equation}\label{eq:nine}
 (WQ)(s)=\frac{1-4}{9!}=-\frac1{120960}.
\end{equation}
This is an instance of the arbitrary-size theorem, not an inference from enumeration.

\paragraph{Completion of the sharp boundary.}
If $n\equiv1\pmod{2w}$ with $n\ge3$, $w\ge2$, write $n=2kw+1$ with $k\ge1$. At $k=1$, Theorem~\ref{thm:balanced} with $\theta=2$ violates canonical stationarity; at $k\ge2$, Theorem~\ref{thm:unit} does so with unit rates. Both use strictly positive OI allocations on the actual recurrent class. Together with Corollary~\ref{cor:positive}, this proves Theorem~\ref{thm:sharp}. In particular, class asymmetry is not necessary for canonical failure on cycles.

\subsection{The five-job four-predecessor calculation}\label{sec:five}
For $w=2$, $s_2=((0),(2,1,3,4))$. At rates $(2,1,1,1,1)$ its canonical weight is $1/48$ and its total outgoing rate is $6$, giving outgoing weighted flow $1/8$. The complete predecessor table, with one-indexed positions, is
\begin{center}
\small
\begin{tabular}{llccc}
\toprule
Previous $c$ & Previous $d$ & Completion & $W_2$ & Rate to $s_2$\\
\midrule
$(0,4)$ & $(2,1,3)$ & $c$: 1 or 2 & $1/36$ & $3$\\
$\emptyq$ & $(2,1,0,3,4)$ & $d$: 1 & $1/240$ & $1$\\
$\emptyq$ & $(2,1,3,4,0)$ & $d$: 3, 4, or 5 & $1/144$ & $4$\\
$\emptyq$ & $(2,3,1,4,0)$ & $d$: 2 & $1/144$ & $1$\\
\bottomrule
\end{tabular}
\end{center}
For a completion in $c$, class $4$ must join the other tail and the prior $c$ contains just $0,4$. For a completion in $d$, class $0$ must join an empty $c$; reversing zero, one, or two replacements gives the other three rows. Applying the rule checks all seven events. Thus
\[
 \frac{3}{36}+\frac{1}{240}+\frac{4}{144}+\frac{1}{144}-\frac{6}{48}
 =\frac{11}{90}-\frac18=-\frac1{360}.
\]
For general $\theta$ the defect reduces to
\[
 (\W Q_{2,\theta})(s_2)
 =-\frac{(\theta-1)(\theta+6)}{24(\theta+2)(\theta+3)(\theta+4)}.
\]
The finite checker derives the complete predecessor list from all 720 configurations, rather than assuming the table is exhaustive.

\section{Excluding every queue-wise factorization at five jobs}\label{sec:certificate}
Theorem~\ref{thm:balanced} excludes the prescribed canonical weights, not an unrelated product $KA(c)B(d)$. We retain a stronger, exact finite result for $w=2,\theta=2$.

\begin{proposition}[Certified nonfactorization]\label{prop:nonfactor}
At rates $(2,1,1,1,1)$ in both queues, the unique stationary law on $\mathcal B_2$ cannot be written as $KA(c)B(d)$, even when $A$ and $B$ are arbitrary queue-local functions.
\end{proposition}
\begin{proof}[Exact certificate]
The file \texttt{data/stationary\_certificate.json} supplies a positive integer weight $v(s)$ for every one of the 180 states of $\mathcal B_2$. Direct integer verification gives
\[
 vQ_{2,2}=0,\qquad
 \sum_{s\in\mathcal B_2}v(s)=29781619500363599243784.
\]
The verifier checks the support, positivity, and every balance equation after reconstructing all transitions. By irreducibility, the stationary law is $v/\sum v$.

Let $c_1=(0,1)$, $c_2=(1,0)$, $d_1=(3,2,4)$, and $d_2=(3,4,2)$. All four pairs $(c_i,d_j)$ belong to $\mathcal B_2$. Their integer weights are
\[
 \begin{pmatrix}
 177097849418281817550&177080689708261370580\\
 506774969345262400380&507080673408934957320
 \end{pmatrix}
 \equiv
 \begin{pmatrix}21&78\\68&29\end{pmatrix}\pmod{101}.
\]
The determinant is congruent to $21\cdot29-78\cdot68\equiv52\pmod{101}$, so it is nonzero. Any factorization $KA(c)B(d)$ would make this rectangle's determinant zero. Normalization does not affect that property.
\end{proof}

This proposition uses an exact computational certificate; the arbitrary-$w$ proof in Section~\ref{sec:balance} does not. An optional rational Gaussian-elimination script regenerates the weights using the 45 orbits under reflection fixing class~0 and queue exchange. Verification is performed separately on all 180 states and does not rely on trusting the solver.


\section{Markov-modulated swap limits}\label{sec:modulated}
Conjecture~2 of~\cite[p.~245]{DG} proposes product form when an exogenous Markov process controls the swap limit. The exceptional class permits a direct test with both limited and unlimited service. Keep $C_5$, one job per vertex, and the additive class rates $(2,1,1,1,1)$ in both queues. Independently of queue events, let the environment jump $L\to U$ and $U\to L$ at rate one. Mode $L$ uses budget two, and mode $U$ allows unlimited replacements; switching changes only the mode. On
\[
 \mathcal S=\mathcal B_2\times\{L,U\},
\]
the joint generator, with the limited states listed first, is
\begin{equation}\label{eq:modulated-generator}
 G=\begin{pmatrix}Q_{2,2}-I&I\\I&Q_{\infty,2}-I\end{pmatrix},
\end{equation}
where $I$ is the $180\times180$ identity matrix. Both queue generators use the same class rates. Each joint state has five queue-completion events of total rate six and one environment transition of rate one.

\begin{theorem}[Nonfactorization under Markov modulation]\label{thm:modulated}
The set $\mathcal S$ is a closed communicating class of 360 states. Its unique stationary law cannot be written as
\begin{equation}\label{eq:modulated-factors}
 \pi(c,d,b)=K(b)A_b(c)B_b(d),
\end{equation}
even with arbitrary real-valued queue factors and prefactors depending on the environment state. In particular, Markov modulation of the swap limit does not guarantee the P\&S product form asserted in Conjecture~2.
\end{theorem}
\begin{proof}
Limited completions preserve $\mathcal B_2$ by Lemma~\ref{lem:closure}; unlimited completions preserve each orientation and hence their union $\mathcal B_2$. Environment jumps preserve the queues. To communicate, switch to mode $L$ if necessary, follow a limited-mode path between the required queue configurations, and switch to the desired final mode. Every step has positive rate.

The file \texttt{data/modulated\_certificate.json} gives a positive integer weight $u(s,b)$ at every joint state. Direct reconstruction of all 2,160 events verifies
\begin{align*}
 uG&=0,\\
 N:=\sum_{(s,b)\in\mathcal S}u(s,b)
   &=453455693620316895843696276658701576451041520896.
\end{align*}
Equivalently, every incoming integer weighted rate equals $7u(s,b)$. Both mode totals equal $N/2$, agreeing with the exogenous environment's stationary law. Irreducibility makes $u/N$ the unique stationary distribution.

Use the same queues $c_1,c_2,d_1,d_2$ as in Proposition~\ref{prop:nonfactor}. At mode $L$, their certified integer weights satisfy
\[
 \begin{pmatrix}
 u(c_1,d_1,L)&u(c_1,d_2,L)\\
 u(c_2,d_1,L)&u(c_2,d_2,L)
 \end{pmatrix}
 \equiv\begin{pmatrix}35&59\\6&1\end{pmatrix}\pmod{101}.
\]
The determinant is congruent to $35\cdot1-59\cdot6\equiv85\pmod{101}$, so it is nonzero. Every factorization~\eqref{eq:modulated-factors} makes this fixed-mode determinant zero, a contradiction. Normalization preserves nonvanishing.
\end{proof}

The graph is three-partite and both allocations are strictly positive and OI. The example starts in the height-three class, which both modes preserve. It concerns modulation of the \emph{budget}; the swapping-graph modulation theorem of~\cite[Theorem~6]{DG} is unaffected. The rational solver uses 90 symmetry orbits, but the independent verifier and Lean certificate check operate on all 360 original states. The Lean development also proves closure, communication, normalization, the generator identity, uniqueness, and exclusion of mode-dependent queue factors.

\section{A queue-length consequence}\label{sec:length}
At unit position rates, $K=|c|$ decreases at rate $K$ and increases at rate $n-K$, independently of job order, swapping graph, and budget. Its birth--death generator therefore determines the same length-process law for any two initial laws with the same length marginal. Every placement word in a closed class occurs with all $n+1$ cuts, so the number of configurations at each length is the same. Since $W(c,d)=1/(|c|!|d|!)$, the normalized canonical length marginal at $j$ is $2^{-n}\binom nj$, the stationary law of this birth--death chain. On $\T$ with $w\ge1$, $k\ge2$, it consequently agrees with the marginal of an actual stationary law, although the canonical candidate is not stationary by Theorem~\ref{thm:unit}. The entire continuous-time length trajectories have identical laws under these two initializations. The nonzero residual $D=\widehat\pi Q$ also satisfies $\sum_{s:|c|=j}D(s)=0$ for every $j$. Detailed path-measure conventions and extensions to length-dependent service capacities are recorded in \texttt{LEAN.md}. This consequence specifies no convergence timescale.

\Needspace{12\baselineskip}
\section{Exhaustive finite results}\label{sec:finite}
The following exact computations supplement the general classification with state counts, head-only service supports, and noncycle graphs.

\subsection{Complete cycle state spaces}
Terminal-component enumeration of all $(n+1)n!$ configurations on $C_5/w=2$ and $C_7/w=3$, for both head-only and all-position completions, gives the following counts. In each graph, $\B$ is the unique tall recurrent class.
\begin{center}
\begin{tabular}{crrrr}
\toprule
$w$ & Full states & Short recurrent & Tall recurrent $=|\B|$ & Transient\\
\midrule
2 & 720 & 480 & 180 & 60\\
3 & 40,320 & 36,400 & 1,120 & 2,800\\
\bottomrule
\end{tabular}
\end{center}
For $w=3$, the 3,248 height-four states split into 1,120 recurrent and 2,128 transient states; all 560 height-five and 112 height-six states are transient. Component sizes are recorded in \texttt{results/classification.json}. Theorem~\ref{thm:classification} gives the all-position structure; these counts and the C7 head-only classification are separately computed.

For the seven nonexceptional pairs $(n,w)\in(\{6,8\}\times\{1,2,3\})\cup\{(7,2)\}$, both supports give no tall recurrent states. Head-only service on $C_6/w=2$ has 68 closed classes, versus 20 for all-position service (\texttt{results/head\_cycles.json}). These are finite checks.

\subsection{Small-graph screen}
With one job per vertex, all 61 nonisomorphic simple bipartite graphs through six vertices pass the $w=1$ screen, including disconnected graphs. Their counts are $1,2,3,7,13,35$; vertex addition and exact canonical labeling verify completeness. Every configuration has a head-only path to the short region, which is closed under all position completions, ruling out tall recurrence.

The same tests at $w=2$ pass for $K_3$, $K_{1,1,2}$, $K_{1,2,2}$, and $C_5$ with one pendant vertex. These 65 graphs cover 192,590 states and 1,143,194 all-position events. The screen settles neither the general one-swap problem nor counterexample minimality.

\subsection{Uniform-formula checks}
All position completions are also checked inside the balanced regions below, at $\theta=1/2,1,2,3$.
\begin{center}
\begin{tabular}{crrr}
\toprule
$w$ & Cycle vertices & $|\B|$ & Position events\\
\midrule
2&5&180&900\\
3&7&1,120&7,840\\
4&9&6,300&56,700\\
5&11&33,264&365,904\\
\bottomrule
\end{tabular}
\end{center}
Two transition implementations check closure, communication within and between orientations, unlimited balance, and the two-flow identity. Complete-space predecessor audits cover $w=2,3$; explicit source words and orientation moves are checked through $w=50$. On the family $n=2w+1$ at $\theta=1$, full canonical balance passes for $w=2,3,4,5$: computational evidence, not an all-$w$ stationarity theorem.


\subsection{Complete exceptional-class check on \texorpdfstring{$C_9$}{C9}}
Two transition implementations check all 131,040 states in $\mathcal T_{9,2}$, its 18 orientations, and 1,179,360 events per generator. With integer weights $\binom9{|c|}=9!W(c,d)$, unlimited residuals all vanish; limited residuals are nonzero at 16,560 states, equal $-3$ at the target in Section~\ref{sec:nine}, and sum to zero at each length. Closure and complete gained/lost predecessor lists are also checked. The counts in \texttt{results/nine\_job.json} are independent checks of Theorem~\ref{thm:unit}.

\section{Conclusion}\label{sec:conclusion}
On cycles, a finite swap budget has an exact structural effect: the recurrent classes are the short orientation fibers and, precisely at $n=2kw+1$, one exceptional class with a single extra edge among otherwise full runs. The operational-to-orientation equivalence makes this a classification of the original queues. It yields a sharp universal canonical boundary for $w\ge2$, with different counterexample allocations covering the first and larger exceptional cycles. At $w=1$, the classification confirms the positive-position conjecture on every even cycle and identifies its two alternating classes. The five-job certificates exclude all queue-wise factors both at a fixed budget and under independent switching between limited and unlimited swaps, supplying a counterexample to Conjecture~2. The head-only argument retains its stronger rate robustness for the fixed-budget class. These finite strengthenings are not asserted for every exceptional cycle.

The classification assumes distinct jobs on a simple cycle and positive service at every position. It does not classify arbitrary swapping graphs, repeated classes, or general supports with vanishing non-head rates. The one-swap problem on general bipartite graphs, removal of positivity from the nonexceptional validity direction, allocation-specific stationarity on exceptional classes, and quantitative convergence rates remain directions for future work.

\section{Verification and reproducibility}\label{sec:verification}
The package includes editable \texttt{paper.tex}, the PDF, every Lean source, exact stationary and orbit certificates, graph inputs, expected results, and regeneration scripts. Running \texttt{python3 run\_checks.py} reconstructs the operational model and reproduces all recorded finite results, also under \texttt{-O}. Both stationary certificates can be regenerated by rational elimination and checked independently against the full generators. The head-only verifier checks all 72,000 event/symmetry identities on the 720 configurations, the paths in Appendix~\ref{app:head}, and forward/reverse head reachability.

Lean verifies Theorems~\ref{thm:classification}, \ref{thm:sharp}, \ref{thm:balanced}, \ref{thm:unit}, and~\ref{thm:modulated}, the C9 specialization, both five-job certificates, and the continuous-time recurrence and queue-length assertions. It constructs the queue transitions, proves the quotient equivalence, and constructs nonexplosive measurable paths from exponential clocks. No external queueing theorem is assumed as an axiom. \texttt{LEAN.md} maps statements to declarations and separates finite computations from general proofs. \texttt{verification/REPORT.md} records the checked revision, commands, and page inspection; \texttt{PACKAGE.md} gives replay and file-hash instructions. CI builds the full library, audits its axioms, replays proofs in the kernel, and runs the Python checks. Only Lean's standard logical axioms are used; no admitted proofs, custom axioms, or native-evaluation axioms are introduced.

\paragraph{Acknowledgment.} This work was assisted by OpenAI's ChatGPT-6 and Anthropic's Claude Fable-5.1.

\Needspace{10\baselineskip}
\begin{thebibliography}{2}
\bibitem{DG} Jan-Pieter Dorsman and Kristen Gardner.
\newblock New directions in pass-and-swap queues.
\newblock \emph{Queueing Systems} \textbf{107}, 205--256 (2024).
\newblock \href{https://doi.org/10.1007/s11134-024-09914-1}{doi:10.1007/s11134-024-09914-1}.
\bibitem{CD} C{\'e}line Comte and Jan-Pieter Dorsman.
\newblock Pass-and-swap queues.
\newblock \emph{Queueing Systems} \textbf{98}, 275--331 (2021).
\newblock \href{https://doi.org/10.1007/s11134-021-09700-3}{doi:10.1007/s11134-021-09700-3}.
\end{thebibliography}

\appendix
\section{The five-job class under arbitrary OI allocation}\label{app:head}
Admissible OI allocations have nonnegative prefix increments and positive singleton capacities~\cite[Definition~1, p.~210]{DG}: head rates are positive, while other rates may vanish.

The 20 symmetries in $\Gamma$ combine $x\mapsto\pm x+k\pmod5$ with optional queue exchange, preserving allowed events but not necessarily rates. Table~\ref{tab:orbits} lists representatives $t_i$ and destination orbits for successive queue positions.
\begin{table}[ht]
\centering\small
\renewcommand{\arraystretch}{1.12}
\begin{tabular}{cllcc}
\toprule
$i$&$c$&$d$&$c$: positions $1,2,\ldots$&$d$: positions $1,2,\ldots$\\
\midrule
1&$\emptyq$&$(0,1,2,4,3)$&--&$(4,5,5,6,6)$\\
2&$\emptyq$&$(0,1,4,2,3)$&--&$(5,4,6,5,5)$\\
3&$\emptyq$&$(0,1,4,3,2)$&--&$(6,6,4,4,4)$\\
4&$(0)$&$(2,1,3,4)$&$(3)$&$(7,7,9,9)$\\
5&$(0)$&$(2,3,1,4)$&$(2)$&$(9,9,7,8)$\\
6&$(0)$&$(2,3,4,1)$&$(1)$&$(8,8,8,7)$\\
7&$(0,1)$&$(2,3,4)$&$(6,6)$&$(9,9,9)$\\
8&$(0,1)$&$(3,2,4)$&$(5,5)$&$(7,7,8)$\\
9&$(0,1)$&$(3,4,2)$&$(4,4)$&$(8,8,7)$\\
\bottomrule
\end{tabular}
\caption{All 45 representative position completions at $w=2$.}\label{tab:orbits}
\end{table}

\begin{proposition}\label{prop:heads}
For every admissible OI allocation, $\mathcal B_2$ is a single closed communicating class of 180 states. At equal unit head-only rates, its uniform distribution is stationary.
\end{proposition}
\begin{proof}
Unequal queue lengths and distinct labels make the action free. Shorter-queue length separates rows 1--3, 4--6, and 7--9; within each group an equivalence fixes two labels and is the identity. Thus the nine balanced orbits exhaust the $180$ states of $\mathcal B_2$.

\Needspace{8\baselineskip}
Closure holds for all positions. Write $C,D$ for head completions. From $t_7$,
\[
 CCDD:t_7\longmapsto((4,3),(2,1,0))=a(t_7),\qquad
 CDD:t_7\longmapsto((0,4,3),(2,1))=b(t_7),
\]
where $a(x)=4-x$ and $b$ combines $x\mapsto2-x$ with queue exchange. They generate $\Gamma$: their product has order ten and $a$ is not one of its powers. Equivariance therefore proves communication within the orbit of $t_7$. The head walk
\[
 7\to6\to1\to4\to3\to6\to8\to5\to2\to5\to9\to4\to7
\]
visits all nine orbits; lifting it through that communicating orbit proves communication throughout $\mathcal B_2$. Every admissible allocation retains these head events with positive rates.

At unit head rates, the table and free group action give one incoming event per state in orbits 1--3 and two in each other orbit. These match the outgoing counts, proving uniform balance for this allocation.
\end{proof}
\end{document}
